\section{Candidate Design Implications for Embodied Neuromorphic Controllers}
\label{sec:implications}
Three of the five implications concern where redundancy pays. Because $50\%$ weight variability and $10\%$ unit-output silencing leave the median command rate unchanged (99.4\% and 100.4\% of baseline), precision and per-unit yield can be traded against wiring redundancy: weights may drift by a lognormal $\sigma$ of 0.5---beyond most analog or multi-bit resistive devices---and 10\% of cores may fail outright, provided many routes converge\cite{arfa2025hardwareaware,thieu2022effects}. Disconnection is paid in latency, not rate---$76.5\%$ of the rate survives at $50\%$ loss while latency rises $33\%$---so a real-time loop should budget that curve (Fig.~\ref{fig:curves}a).

The other two implications concern failure modes that per-unit monitoring misses. Hub silencing is the only class that can \emph{amplify}, and its sign matters more than its size: removing the inhibitory half of the top 1\% raises activity $21\times$, while the full set hides this behind a matching excitatory loss---so monitor the activity envelope, not liveness; and since 8.6\% of inhibitory mass removed at random was benign in eleven of twelve trials while 5.6\% removed at hubs always destabilised, sparing high-fanout \emph{inhibitory} units is the cheaper mitigation. Deafferentation is the opposite extreme: either channel alone retains 68--79\% of the median rate, but losing both drives the pathway to zero, making the sensor interface---not the synaptic fabric---the first redundancy candidate. Because relative pathway decay is identical while only the biological graph yields a decodable command, fault injection belongs pre-deployment\cite{yik2023neurobench,medvedeva2025formalizing}.

\enlargethispage{3\baselineskip}
\section{Limitations and Conclusion}
Our LIF abstraction fixes one drive rate, the hub sign split is measured at that one rate, and correlated damage is excluded. One pathway is instrumented, so the ordering is a hypothesis for others rather than a whole-brain law, and its vulnerability to aging\cite{gaitanidis2025drosophila} suggests tolerance is not fixed either. Under the tested drive the median rate survives weight noise and unit silencing, retaining 76.5\% at 50\% loss; the dangerous failure is inhibitory-hub over-activity, which a closed loop turns into lost steering.
